
\begin{frame}{Le procédé de photolithographie}

\begin{columns}[T,onlytextwidth]

  % ---- Colonne de gauche : points clés ----
  \begin{column}{0.47\textwidth}
    \footnotesize
    \begin{itemize}\setlength{\itemsep}{3pt}
      \item \textbf{Principe} : transférer un motif géométrique d'un masque vers une résine photosensible déposée sur une plaquette (\emph{wafer})
      \item \textbf{Dépôt} : enduction centrifuge (\emph{spin coating}) d'un film uniforme
      \item \textbf{Insolation} : exposition UV à travers un masque optique portant le motif
      \item \textbf{Développement} : dissolution sélective des zones exposées (résine positive) ou non exposées (négative)
      \item \textbf{Gravure} : transfert du motif dans le substrat (voie humide ou plasma)
      \item \textbf{Retrait} de la résine (\emph{stripping})
      \item Résolution limitée par la diffraction $\Rightarrow$ UV profond, immersion, EUV
    \end{itemize}
  \end{column}

  % ---- Colonne de droite : schéma ----
  \begin{column}{0.51\textwidth}
    \centering
    \includegraphics[width=\textwidth]{./images/photolitho_v2.pdf}
  \end{column}

\end{columns}

\end{frame}
